Euler's formula $V-E+F=2$ would give $V=2+18-12=8$. Every edge has two ends, so the sum of vertex degrees would be $2E=36$. In a regular polyhedron every vertex has the same integer degree $q$, requiring $8q=36$, or $q=\dfrac{9}{2}$, impossible. Thus no such regular polyhedron exists, regardless of an asserted edge length, and it has no defined volume.
\begin{asy}
texpreamble("\usepackage{amsmath,amssymb}");
size(240);
defaultpen(fontsize(9pt));
pair[] v={(0,0),(3,0),(3,3),(0,3),(1,1),(4,1),(4,4),(1,4)};int[][] e={{0,1},{1,2},{2,3},{3,0},{4,5},{5,6},{6,7},{7,4},{0,4},{1,5},{2,6},{3,7}};for(int[] z:e)draw(v[z[0]]--v[z[1]],(z[0]==4||z[1]==4)?dashed:solid);for(int i=0;i<8;++i){dot(v[i]);label(string(i+1),v[i],(i==4?(-1,0):(v[i]-(2,2))));}label("Example only: $V=8,\ E=12,\ F=6$",(2,-.8));
\end{asy}
